\chapter{Methodology}

\section{Experimental Environment}

All benchmarks run in the RISC Zero host/guest model. Guest programs execute
cryptographic logic inside the zkVM, while host programs serialize inputs,
trigger proving, and verify receipts.

RISC Zero is used because it gives a concrete and auditable contract for this
project: a guest image ID binds the proved program, the journal exposes public
outputs, and receipt verification gives a uniform success condition across all
targets \cite{risczeroDocs}. Alternative zkVMs are relevant, but including them
would have changed the project from a controlled single-stack study into a
platform-comparison study. That broader design is therefore left to future work.

The harness records machine-readable artifacts per run, including target command
metadata, per-trial measurements, and success/failure status for each trial.
Core outputs are aggregated into \texttt{aggregated.json} files for both
primitive and operation-level campaigns.

The report benchmark manifests record the run ID, target commands, Python
runtime, and platform tuple. For the main reported campaigns, this provenance is
Darwin 25.3.0 on arm64 with Python 3.9.6; the repository lockfiles pin
\texttt{risc0-zkvm} 3.0.5 for the final AES targets. This is useful provenance,
but it is not a complete hardware identity record, so timing results are still
treated as environment-specific.

This chapter treats the harness as part of the scientific method: every major
claim in later chapters maps to a benchmark family and fixed target
configuration.

\section{Metrics}

The evaluation uses zkVM-relevant metrics rather than native CPU timing alone:

\begin{itemize}
\item proving time and verification time;
\item total, user, paging, and reserved cycle counts;
\item success rate across repeated trials;
\item proof-size fields (\texttt{proof\_bytes} and
\texttt{full\_receipt\_bytes}) where available.
\end{itemize}

The primary optimization target is proving time, with cycle and proof-size
metrics used to explain observed behavior.

Cycle metrics are interpreted with distinct roles. \(C_{user}\) reflects guest
algorithm work, so it is used as the algorithm-efficiency signal. \(C_{total}\)
reflects the trace footprint that is actually proven, so it is used to explain
proof-system load boundaries. \(t_{prove}\) captures the prover's wall-clock
cost, which is the practical resource cost for evidence generation.

This distinction matters because \(C_{total}\) is constrained by the proof
system, not only by guest instruction count. RISC Zero documents cycles and
execution traces as the core cost surface, while the STARK/FRI proof structure
uses polynomial evaluation domains \cite{risczeroTraceDocs,bensasson2018stark}.
Chapter 11 therefore treats the power-of-two \(C_{total}\) buckets observed in
the benchmark data as an empirical trace-footprint effect rather than as native
CPU timing behavior.

These metrics are appropriate because they correspond to different stakeholder
questions. Proving time measures the cost paid by the party generating evidence,
verification time measures the cost paid by the checker, cycle counts provide a
more execution-structure-sensitive signal than wall-clock time alone, and proof
or receipt sizes affect storage and transport overhead. Reporting all of them
reduces the risk of optimizing for a single number that hides a deployment cost.

Medians are the headline statistic because prover measurements can be
right-skewed by occasional slow runs. Where only one trial is available, the
reported median is simply the observed value and is treated as directional rather
than as a high-power estimate.

Native CPU benchmarks are not used as headline evidence because the research
question is proof cost, not conventional execution speed. Similarly,
operation-level rows are used only to interpret likely hotspots; full-cipher and
full-pipeline targets remain authoritative because they include serialization,
guest execution, receipt generation, and verification.

\section{Benchmark Scope}

The locked benchmark scope consists of:

\begin{itemize}
\item operation-level suite: \texttt{operation-bnchmrk-r0};
\item baseline ciphers: \texttt{aes-r0}, \texttt{salsa-r0}, and
\texttt{lowmc-r0};
\item optimization variants: \texttt{aes-r0-optimised} and
\texttt{lowmc-r0-optimised};
\item final authenticated pipeline split into two benchmark families:
CTR-only block-sweep targets (for example \texttt{aes-ctr-1blk},
\texttt{aes-ctr-4blk}, \texttt{aes-ctr-16blk}, \texttt{aes-ctr-64blk}) and
CTR+HMAC compare targets (for example \texttt{aes-ctr-hmac-4blk}).
\end{itemize}

This split is structural in the repository and harness configuration: CTR-only
and CTR+HMAC are separate target families (\texttt{aes-ctr} and
\texttt{aes-ctr-hmac}), rather than one target controlled by a runtime mode
flag. The distinction improves traceability because benchmark IDs encode both
mode family and workload size.

Operation benchmarks are treated as directional evidence for optimization
priority. They are not used as a standalone ISA-level micro-architectural truth
and are cross-checked against full-cipher runs.

CTR versus CTR+HMAC block-size parity is complete in the report benchmark run
for 1, 4, 16, and 64 blocks, enabling like-for-like overhead interpretation in
Chapter 11.

\section{Trial Count Rationale and Statistical Power}

The report benchmark campaign uses \(N=1\) per target, where \(N\) denotes the
number of independently timed proving trials run for each benchmark target, to
cover the full set of cipher and mode-family comparisons within dissertation
compute/time limits, especially for long-running LowMC proving jobs.

This supports tentative directional ranking for large effect-size comparisons and
evidence-trace demonstration, but not high-confidence inference over small
differences. To avoid overstating precision, the report:

\begin{itemize}
\item labels one-trial values as directional where appropriate;
\item prioritizes effect sizes that are large relative to noise expectations;
\item references earlier multi-trial pilot runs where available.
\end{itemize}

For final production-grade benchmarking, the intended policy is \(N \geq 5\)
for all headline targets, with variability reporting (mean/p95/stddev) and
explicit outlier commentary.

The main alternative design would have been a narrower benchmark matrix with
\(N \geq 5\) immediately. That would improve statistical power, but would remove
coverage of either LowMC/AES comparison, authenticated-mode overhead, or
operation-level context. The chosen design prioritizes breadth and traceability
for dissertation evidence, then treats small deltas cautiously. This is why large
effect-size results, such as the LowMC-versus-AES proving-time gap, carry more
weight than sub-second wall-clock differences in the authentication comparison.

\section{Threats to Validity and Confounds}

Key confounds in this benchmark setting are:

\begin{itemize}
\item \textbf{Thermal and power effects:} long prover runs can trigger clock
adjustments;
\item \textbf{Background load:} OS scheduling and concurrent processes can add
timing jitter;
\item \textbf{Build-state effects:} initial compiles and cache state can distort
first-run timing;
\item \textbf{Code-drift risk:} missing explicit commit embedding in artifacts can
complicate strict reruns. If the guest codebase changes between benchmark
campaigns and those changes are not recorded in the artifact JSON, a later run
may silently compare different implementations. The harness records command and
path provenance but does not currently embed per-target git revisions; this
limitation is noted in Chapter 12 and would be removed by adding a
\texttt{git rev-parse HEAD} field to the artifact schema before each run.
\end{itemize}

Mitigations used in this dissertation are fixed harness configs, explicit run
IDs, machine-readable per-trial artifacts, and clear separation between
directional and final-claim interpretation.
